\chapter{Introduction}
\label{chap:intro}

\section{Background and motivation}

Rapid urbanisation is increasing demand for efficient passenger transport and urban goods delivery. The United Nations projects that approximately 68\% of the world's population will live in urban areas by 2050 \cite{unwup2019}. This growth, together with the need to reduce transport emissions and operating costs, has encouraged the development of lightweight electric-drive vehicles for last-mile logistics. Within this context, the TUKZIE Flagship Programme seeks to engineer an energy-aware, secure and autonomous-ready vehicle platform for last-mile logistics and urban freight applications. The TUKZIE Rev 0, a 72V electric cargo trike, represents a testbed for integrating advanced telemetry and driver-assistance systems in resource-constrained environments.

The relevance of this platform is particularly clear in sub-Saharan African cities, where three-wheelers already support urban goods movement. A 2025 study of medium-size cargo delivery routes in Dar es Salaam found that an electric tricycle reduced operating cost per kilometre by approximately 45.5\% relative to a motorcycle, 62.5\% relative to a petrol tuk-tuk, and 86\% relative to a light-duty car \cite{sebwa2025}. The reported values are route- and price-specific rather than universal forecasts, but they show why locally appropriate electric freight platforms warrant further engineering and field validation.

Vehicle telemetry is often only logged and analysed afterwards, which delays maintenance and prevents real-time decisions on the vehicle. Real-time safety alerts matter because road traffic injuries remain a leading cause of death worldwide, with most of the deaths in low- and middle-income countries \cite{who2023}. Presenting telemetry to the driver in real time has its own cost, however: any driver display, head-down or head-up, draws the driver's gaze or attention away from the road, and cluttered presentation increases attention capture \cite{liu2003, hauslschmid2015}. NHTSA guidance recommends that individual visual-manual glances away from the roadway should not exceed two seconds, with no more than twelve seconds of cumulative off-road glance time for a task \cite{nhtsa2013}. This provides a measurable human-factors constraint on any driver display: what is shown must be sparse and legible enough to be read within a short glance, wherever the display is placed.

Recent advances in High-Performance Embedded Computing (HPEC) have enabled edge-based signal processing in vehicular applications, allowing for real-time sensor fusion and on-vehicle analytics \cite{barnell2018}. Many light vehicles, the TUKZIE Rev 0 among them, already carry a touchscreen dashboard on which such analytics can be presented; Head-Up Displays (HUDs) offer a forward-view alternative \cite{mareska2022, canomarin2016}, with optical and attention costs of their own (Chapter~\ref{chap:lit}). However, integrating on-vehicle analytics with driver presentation for low-cost, three-wheeled electric vehicles in developing contexts remains an open challenge.

This study addresses this gap by developing an HPEC telemetry unit for the TUKZIE Rev 0, integrated with the vehicle's existing dashboard, enabling real-time ride-smoothness characterisation and safety alerts to the driver. The original brief specified a windshield HUD. During the project it was replaced, with the supervisor's agreement, by integration with the existing dashboard: the dashboard already shows the information a HUD would have carried, a second display would compete for the driver's attention \cite{zhu2021}, and the programme's move towards autonomy makes the camera's hazard alerts more useful than a second screen (Section~\ref{sec:dashboard-page}).

The intended outcome is an integrated prototype demonstrated on the target or a representative vehicle under relevant operating conditions. This corresponds to a target of Technology Readiness Level 6, at which a representative prototype is demonstrated in a relevant environment \cite{kimmel2020}. It does not imply production qualification or approval for unrestricted road deployment.

\section{Objectives of this study}

\subsection{Problems to be investigated}

Transportation agencies and fleet operators face significant challenges in monitoring road health and vehicle performance cost-effectively. Reference-grade profiling equipment requires specialist instrumentation and operation, which limits the frequency and coverage of surveys in resource-constrained settings \cite{fhwa2016, alatoom2025}. Standard in-car displays often contribute to cognitive overload and attention capture, where drivers divert their gaze from the road to process telemetry \cite{zhu2021}. Furthermore, in resource-constrained environments, there is a gap for a modular, low-latency system capable of edge-based signal processing for vehicle dynamic-response characterisation that also presents its results safely on the driver's existing display.

The engineering problem is therefore to acquire synchronised data reliably on a high-vibration, power-constrained vehicle; distinguish road-induced motion from vehicle-specific dynamics; maintain bounded processing and display latency; and preserve telemetry during intermittent cellular connectivity. These functions must be integrated without obstructing the driver's view or making safety-critical behaviour dependent on a remote service \cite{winberg2026}.

Specifically, this study investigates the following research questions:

\begin{enumerate}
\item How accurately and repeatably can a low-cost edge-processing system characterise the trike's dynamic response to road- and powertrain-induced excitation in real time, using inertial sensors on a 72V electric cargo trike?
\item Which information hierarchy, layout, and update strategy can present safety-critical telemetry and hazard alerts on the vehicle's existing dashboard while limiting visual clutter and driver distraction?
\item What end-to-end latency, data-loss, and recovery performance can be achieved across sensor acquisition, edge processing, local logging, and 4G/LTE synchronisation?
\item Which time- and frequency-domain features most effectively separate road-induced vibration from repeatable motor and powertrain vibration?
\end{enumerate}

\subsection{Requirements and verification mapping}
\label{sec:requirements-mapping}

Table~\ref{tab:requirements-mapping} maps each research question to a concrete system requirement and the method by which it is, or will be, verified. This mapping is used consistently through Chapter~\ref{chap:method} (Methodology) and is intended to be revisited directly in Chapter~\ref{chap:results}, rather than treated as a fixed statement made only here.

\begin{table}[ht]
\centering
\caption{Research questions mapped to system requirements and verification method.}
\label{tab:requirements-mapping}
\begin{tabular}{|p{0.9cm}|p{5.4cm}|p{5.4cm}|p{2.2cm}|}
\hline
\textbf{RQ} & \textbf{Requirement} & \textbf{Verification method} & \textbf{Status} \\ \hline
RQ1 & Dual-IMU acquisition at $\geq$160\,Hz (ISO~2631-1 bound) with per-sensor magnitude agreement to within measurement noise & Deterministic 200\,Hz acquisition (Section~\ref{sec:dual-imu-architecture}); per-sensor calibration against true gravity, cross-checked under a dynamic disturbance (Section~\ref{sec:imu-magnitude-calibration}) & Bench-verified \\ \hline
RQ2 & The existing dashboard presents safety-critical alerts without exceeding NHTSA off-road glance-time guidance & Sparse alert page in the dashboard's own status colours, tested against the live camera (Section~\ref{sec:dashboard-page}); on-vehicle glance-time observation & Page bench-verified; glance time outstanding \\ \hline
RQ3 & Bounded end-to-end latency, quantified data loss, and recovery across acquisition, edge processing, local logging, and cellular sync & Queue drop-counters and jitter instrumentation (Section~\ref{sec:acquisition-instrumentation}); local LittleFS logging decoupled from MQTT connection state; end-to-end MQTT publish against a real broker (Section~\ref{sec:mqtt-networking}) & Acquisition, camera latency, signal-loss recovery and local-log read-back bench-verified; on-vehicle testing outstanding \\ \hline
RQ4 & Time/frequency-domain features that separate road-induced from powertrain-induced vibration & Closed-form self-test of feature computation against a synthetic signal (Section~\ref{sec:ride-characterisation}); real-vehicle feature comparison against known road/powertrain events & Self-test only; real-vehicle validation outstanding \\ \hline
\end{tabular}
\end{table}

\subsection{Purpose of the study}

The significance of investigating these problems lies in bridging the gap between expensive, proprietary road-profiling equipment and the need for cost-effective, real-time monitoring solutions. By leveraging High-Performance Embedded Computing and low-cost inertial sensors, this study aims to provide a modular, low-latency system that delivers head-up telemetry while limiting unnecessary downward glances and avoiding additional visual clutter.

The findings will directly inform the development of the TUKZIE Rev 0 platform and contribute to the broader field of sustainable urban mobility by demonstrating a replicable, open-architecture approach to vehicle telemetry. This work aligns with the UCT Radar Remote Sensing Group's focus on advanced sensing and the TUKZIE Flagship Programme's mandate for energy-aware, secure vehicle platforms.

Furthermore, this research addresses the practical needs of fleet operators in developing regions, where infrastructure monitoring budgets are constrained and traditional data collection methods are often infeasible. By enabling continuous, vehicle-specific ride-condition screening, the proposed system could support route comparison, targeted inspection, and more informed maintenance planning.

\section{Scope and Limitations}

\subsection{Scope}

This study focuses exclusively on the TUKZIE Rev 0, a 72\,V electric cargo trike manufactured by Jinpeng \cite{jinpengmanual}. Direct inspection of the motor and controller nameplates recorded a three-phase, permanent-magnet, sensored hub motor rated at 72\,V DC and 35\,A, paired with a controller rated to a 90\,A current limit, giving a calculated rated input power of approximately 2.5\,kW \cite{adamsrasara2026}. The vehicle's own manufacturer identification plate (Figure~\ref{fig:vehicle-nameplate}, Appendix~\ref{app:figures}), photographed directly on the platform, independently confirms the manufacturer and model (Jiangsu Jinpeng Group Co., Ltd, model L2e-U) and additionally records a rated output of 3.9\,kW, a rated speed of 42\,km/h, and a maximum load of 639\,kg. The system integrates the following components:


\begin{itemize}
\item A hardware module housing inertial measurement units (IMUs) and a Sensirion SEN55 environmental sensor, which measures particulate matter, volatile organic compounds (VOC), nitrogen oxides (NOx), humidity, and temperature (Section~\ref{sec:environmental-sensing}).
\item A forward-facing camera, time-synchronised to the inertial acquisition path, providing a visual record against which detected dynamic-response events can be labelled and interpreted (for example, distinguishing a pothole from an intersection or a deliberate manoeuvre), and running a pretrained object-detection model to report predefined road-relevant object categories and an approximate distance band (Section~\ref{sec:hazard-detection}).
\item A non-visual, non-machine-learning proximity sensor (an 8$\times$8-zone time-of-flight distance sensor) providing live nearby-object distance and rough left/centre/right position to the driver, adopted as a second, physically independent measurement that operates regardless of the camera detector's own status and involves no camera image or learned model.
\item A multithreaded software architecture for real-time sensor fusion, telemetry processing, and dynamic-response indexing.
\item A camera page for the vehicle's existing dashboard, showing the live view and the active hazard alerts in the dashboard's own status colours, in place of the optical windshield Head-Up Display of the original brief (Section~\ref{sec:dashboard-page}).
\item A 4G/LTE communication link for central database synchronisation.
\end{itemize}

The project is aligned with the UCT Radar Remote Sensing Group's focus on advanced sensing and specifically addresses the ``ride-smoothness'' mandate of TUKZIE by isolating road-induced vibration from nominal motor and powertrain vibration using statistical feature extraction. The scope is limited to the TUKZIE Rev 0 platform and utilises low-cost inertial and environmental sensors processed at the edge using microcontroller and single-board computer platforms \cite{alatoom2025}.

The software architecture follows a modular design, with separate threads for:

\begin{itemize}
\item High-frequency sensor polling (both IMUs at 200\,Hz, above the 160\,Hz minimum in Table~\ref{tab:requirements-mapping}).
\item Compute-heavy signal-processing algorithms (ride characterisation).
\item The driver display (the dashboard page, on the dashboard's own computer).
\item Networking stack (MQTT over 4G LTE Cat-1).
\end{itemize}

This separation ensures that no process blocks another \cite{winberg2026}.

\subsection{Limitations}

The following limitations are acknowledged:

\begin{itemize}
\item The study is limited to a single vehicle platform (TUKZIE Rev 0), which may affect generalisability to other vehicle types or configurations.
\item Sensor calibration assumes nominal operating conditions; significant environmental variations (extreme temperatures and humidity) may affect accuracy.
\item Edge-processing latency measurements are conducted in controlled test environments; real-world performance may vary with network conditions and electromagnetic interference.
\item The driver display is the vehicle's existing head-down dashboard. A windshield head-up display, which would keep the driver's gaze nearer the road \cite{mareska2022}, is not built in this project.
\item Time constraints limit extensive field testing to urban routes only; rural and off-road conditions are not evaluated.
\item The system is designed as a removable module to ensure that the vehicle can be quickly returned to a standard setup without obstructions \cite{winberg2026}.
\item The 4G/LTE communication depends on network availability; data synchronisation may be delayed in areas with poor coverage.
\item The camera's object-detection capability uses a pretrained, general-purpose model without vehicle-specific training, reports only object category, approximate distance band, detection confidence, timestamp, and alert status, and does not retain raw camera frames or identify individual people; it does not track objects across frames or perform lane-level blind-spot geometry, and is not used for driver monitoring or any safety-critical vehicle function. The time-of-flight rangefinder described above provides a second, independent proximity measurement alongside it.
\end{itemize}

\section{Plan of development}

This report is organised as follows:

\textbf{Chapter 2: Literature Review.} Reviews telemetry systems for electric vehicles, edge-computing architectures for vehicular applications, inertial road-roughness assessment \cite{alatoom2025, bui2026}, head-up display design \cite{mareska2022, canomarin2016, hauslschmid2016}, and proximity sensing on low-speed vehicles.

\textbf{Chapter 3: Methodology.} Describes the hardware platform, firmware architecture, sensor integration, calibration, camera and proximity subsystems, and cellular telemetry, together with the bench validation carried out at each stage, including faults found and how they were resolved.

\textbf{Chapter 4: Results.} Presents acquisition integrity, calibration, ride-characterisation, battery, camera, proximity and cellular results by subsystem, followed by a full-system test on the vehicle.

\textbf{Chapter 5: Discussion.} Interprets the results against the literature and the research questions.

\textbf{Chapter 6: Conclusions.} Answers each research question and states the contribution and its limitations.

\textbf{Chapter 7: Recommendations.} Identifies future work, including generic road-obstacle detection.
